\subsection{DISPLACEMENTS}

We recall again that the affine transformations \(f\) of \(\mathbf{R}^n\) onto itself which leave invariant the distance between any two points {[}that is to say, such that \(d(f(x), f(y)) = d(x, y)\) for all \(x, y]\) are called Euclidean displacements (or simply displacements) (*); they form a group, called the group of displacements of \(\mathbf{R}^n\) . This group operates transitively on \(\mathbf{R}^n\) ; more generally, if \(V\) and \(V'\) are any two affine linear varieties of the same dimension in \(\mathbf{R}^n\) , there exists a displacement which transforms \(V\) into \(V'\) . The displacements which leave the origin fixed, called orthogonal transformations, form a subgroup of the group of all displacements. This subgroup is called the orthogonal group on \(n\) real variables; the linear mappings which belong to this group are characterized by the fact that they leave invariant the norm \(||x||\) of every point \(x \in \mathbf{R}^n\) , or, equivalently, the quadratic form \(||x||^2 = \sum_{i=1}^{n} x_i^2\) . The scalar product of two vectors \(x = (x_i)\) and \(y = (y_i)\) of \(\mathbf{R}^n\) is the value \(\sum_{i=1}^{n} x_i y_i\) of the bilinear form associated with the quadratic form \(\frac{1}{2} \sum_{i=1}^{n} x_i^2\) ; it is denoted by \((x|y)\) , or simply by \(xy\) if there is no likelihood of confusion. Every orthogonal transformation leaves invariant the scalar product of any two vectors. Two vectors \(x, y\) are said to be orthogonal if \((x|y) = 0\) ; two vector subspaces \(V, V'\) of \(\mathbf{R}^n\) are said to be orthogonal if each \(x \in V\) is orthogonal to each \(y \in V'\) ; and two affine linear varieties \(P, P'\) are said to be orthogonal if the vector subspaces parallel respectively to \(P\) and \(P'\) are orthogonal.

